% Options for packages loaded elsewhere
\PassOptionsToPackage{unicode}{hyperref}
\PassOptionsToPackage{hyphens}{url}
%
\documentclass[
]{article}
\usepackage{amsmath,amssymb}
\usepackage{lmodern}
\usepackage{iftex}
\ifPDFTeX
  \usepackage[T1]{fontenc}
  \usepackage[utf8]{inputenc}
  \usepackage{textcomp} % provide euro and other symbols
\else % if luatex or xetex
  \usepackage{unicode-math}
  \defaultfontfeatures{Scale=MatchLowercase}
  \defaultfontfeatures[\rmfamily]{Ligatures=TeX,Scale=1}
\fi
% Use upquote if available, for straight quotes in verbatim environments
\IfFileExists{upquote.sty}{\usepackage{upquote}}{}
\IfFileExists{microtype.sty}{% use microtype if available
  \usepackage[]{microtype}
  \UseMicrotypeSet[protrusion]{basicmath} % disable protrusion for tt fonts
}{}
\makeatletter
\@ifundefined{KOMAClassName}{% if non-KOMA class
  \IfFileExists{parskip.sty}{%
    \usepackage{parskip}
  }{% else
    \setlength{\parindent}{0pt}
    \setlength{\parskip}{6pt plus 2pt minus 1pt}}
}{% if KOMA class
  \KOMAoptions{parskip=half}}
\makeatother
\usepackage{xcolor}
\usepackage{longtable,booktabs,array}
\usepackage{multirow}
\usepackage{calc} % for calculating minipage widths
% Correct order of tables after \paragraph or \subparagraph
\usepackage{etoolbox}
\makeatletter
\patchcmd\longtable{\par}{\if@noskipsec\mbox{}\fi\par}{}{}
\makeatother
% Allow footnotes in longtable head/foot
\IfFileExists{footnotehyper.sty}{\usepackage{footnotehyper}}{\usepackage{footnote}}
\makesavenoteenv{longtable}
\usepackage{graphicx}
\makeatletter
\def\maxwidth{\ifdim\Gin@nat@width>\linewidth\linewidth\else\Gin@nat@width\fi}
\def\maxheight{\ifdim\Gin@nat@height>\textheight\textheight\else\Gin@nat@height\fi}
\makeatother
% Scale images if necessary, so that they will not overflow the page
% margins by default, and it is still possible to overwrite the defaults
% using explicit options in \includegraphics[width, height, ...]{}
\setkeys{Gin}{width=\maxwidth,height=\maxheight,keepaspectratio}
% Set default figure placement to htbp
\makeatletter
\def\fps@figure{htbp}
\makeatother
\setlength{\emergencystretch}{3em} % prevent overfull lines
\providecommand{\tightlist}{%
  \setlength{\itemsep}{0pt}\setlength{\parskip}{0pt}}
\setcounter{secnumdepth}{-\maxdimen} % remove section numbering
\ifLuaTeX
  \usepackage{selnolig}  % disable illegal ligatures
\fi
\IfFileExists{bookmark.sty}{\usepackage{bookmark}}{\usepackage{hyperref}}
\IfFileExists{xurl.sty}{\usepackage{xurl}}{} % add URL line breaks if available
\urlstyle{same} % disable monospaced font for URLs
\hypersetup{
  hidelinks,
  pdfcreator={LaTeX via pandoc}}

\author{}
\date{}

\begin{document}

\begin{longtable}[]{@{}
  >{\raggedright\arraybackslash}p{(\columnwidth - 2\tabcolsep) * \real{0.5000}}
  >{\raggedright\arraybackslash}p{(\columnwidth - 2\tabcolsep) * \real{0.5000}}@{}}
\toprule()
\begin{minipage}[b]{\linewidth}\raggedright
\begin{quote}
ICS1512
\end{quote}
\end{minipage} & \begin{minipage}[b]{\linewidth}\raggedright
Machine Learning Algorithms Laboratory
\end{minipage} \\
\midrule()
\endhead
\bottomrule()
\end{longtable}

\begin{longtable}[]{@{}
  >{\raggedright\arraybackslash}p{(\columnwidth - 2\tabcolsep) * \real{0.5000}}
  >{\raggedright\arraybackslash}p{(\columnwidth - 2\tabcolsep) * \real{0.5000}}@{}}
\toprule()
\begin{minipage}[b]{\linewidth}\raggedright
\begin{quote}
Experiment No.
\end{quote}
\end{minipage} & \begin{minipage}[b]{\linewidth}\raggedright
\begin{quote}
02
\end{quote}
\end{minipage} \\
\midrule()
\endhead
\begin{minipage}[t]{\linewidth}\raggedright
\begin{quote}
Experiment Title
\end{quote}
\end{minipage} & Email Spam/Ham Classification using Naive Bayes and KNN \\
\textbf{GitHub Repository:} & \url{https://github.com/The-FOOL-00/ML_Lab} \\
\begin{minipage}[t]{\linewidth}\raggedright
\begin{quote}
Student Name
\end{quote}
\end{minipage} & \begin{minipage}[t]{\linewidth}\raggedright
\begin{quote}
R. Sundareswaran
\end{quote}
\end{minipage} \\
\begin{minipage}[t]{\linewidth}\raggedright
\begin{quote}
Register Number
\end{quote}
\end{minipage} & \begin{minipage}[t]{\linewidth}\raggedright
\begin{quote}
3122247001048
\end{quote}
\end{minipage} \\
\begin{minipage}[t]{\linewidth}\raggedright
\begin{quote}
Faculty
\end{quote}
\end{minipage} & \begin{minipage}[t]{\linewidth}\raggedright
\begin{quote}
Dr. Poreddy Ajaykumar Reddy
\end{quote}
\end{minipage} \\
\begin{minipage}[t]{\linewidth}\raggedright
\begin{quote}
Submission Date
\end{quote}
\end{minipage} & \begin{minipage}[t]{\linewidth}\raggedright
\begin{quote}
August 11, 2026
\end{quote}
\end{minipage} \\
\bottomrule()
\end{longtable}

\begin{longtable}[]{@{}
  >{\raggedright\arraybackslash}p{(\columnwidth - 2\tabcolsep) * \real{0.5000}}
  >{\raggedright\arraybackslash}p{(\columnwidth - 2\tabcolsep) * \real{0.5000}}@{}}
\toprule()
\begin{minipage}[b]{\linewidth}\raggedright
\begin{quote}
\textbf{1}
\end{quote}
\end{minipage} & \begin{minipage}[b]{\linewidth}\raggedright
\begin{quote}
\textbf{Objective}
\end{quote}
\end{minipage} \\
\midrule()
\endhead
\bottomrule()
\end{longtable}

\begin{quote}
• Build Spam classifier using Naive Bayes and KNN.

• Compare Gaussian, Multinomial and Bernoulli Naive Bayes.• Study the
effect of varying k.

• Compare KDTree and BallTree.

• Optimize KNN using GridSearchCV and RandomizedSearchCV.• Evaluate
computational complexity and execution time.

• Evaluate using K-Fold Validation.

\textbf{2} \textbf{Problem Statement}\\
The objective of this experiment is to develop an Email Spam/Ham
Classification model using supervised machine learning algorithms. The
experiment compares the performance of three variants of Naive Bayes
classifiers (Gaussian, Multinomial and Bernoulli) with the K-Nearest
Neighbors (KNN) algorithm. Hyperparameter tuning is performed using
GridSearchCV and RandomizedSearchCV to identify the optimal KNN
configuration. The performance of KDTree and BallTree search algorithms
is also analyzed along with computational complexity, execution time and
K-Fold Cross Validation.
\end{quote}

\begin{longtable}[]{@{}
  >{\raggedright\arraybackslash}p{(\columnwidth - 2\tabcolsep) * \real{0.5000}}
  >{\raggedright\arraybackslash}p{(\columnwidth - 2\tabcolsep) * \real{0.5000}}@{}}
\toprule()
\begin{minipage}[b]{\linewidth}\raggedright
\begin{quote}
\textbf{3}
\end{quote}
\end{minipage} & \begin{minipage}[b]{\linewidth}\raggedright
\begin{quote}
\textbf{Dataset Description}
\end{quote}
\end{minipage} \\
\midrule()
\endhead
\bottomrule()
\end{longtable}

\begin{longtable}[]{@{}
  >{\raggedright\arraybackslash}p{(\columnwidth - 2\tabcolsep) * \real{0.5000}}
  >{\raggedright\arraybackslash}p{(\columnwidth - 2\tabcolsep) * \real{0.5000}}@{}}
\toprule()
\begin{minipage}[b]{\linewidth}\raggedright
\begin{quote}
Dataset Name
\end{quote}
\end{minipage} & \begin{minipage}[b]{\linewidth}\raggedright
\begin{quote}
Spambase
\end{quote}
\end{minipage} \\
\midrule()
\endhead
\begin{minipage}[t]{\linewidth}\raggedright
\begin{quote}
Dataset Source
\end{quote}
\end{minipage} & \begin{minipage}[t]{\linewidth}\raggedright
\url{https://www.kaggle.com/datasets/somesh24/spambase}
\end{minipage} \\
\begin{minipage}[t]{\linewidth}\raggedright
\begin{quote}
Number of Samples
\end{quote}
\end{minipage} & \begin{minipage}[t]{\linewidth}\raggedright
\begin{quote}
4601
\end{quote}
\end{minipage} \\
\begin{minipage}[t]{\linewidth}\raggedright
\begin{quote}
Number of Features
\end{quote}
\end{minipage} & \begin{minipage}[t]{\linewidth}\raggedright
\begin{quote}
58
\end{quote}
\end{minipage} \\
\begin{minipage}[t]{\linewidth}\raggedright
\begin{quote}
Number of Classes
\end{quote}
\end{minipage} & \begin{minipage}[t]{\linewidth}\raggedright
\begin{quote}
2
\end{quote}
\end{minipage} \\
\begin{minipage}[t]{\linewidth}\raggedright
\begin{quote}
Missing Values
\end{quote}
\end{minipage} & \begin{minipage}[t]{\linewidth}\raggedright
\begin{quote}
0
\end{quote}
\end{minipage} \\
\begin{minipage}[t]{\linewidth}\raggedright
\begin{quote}
Train-Test Split
\end{quote}
\end{minipage} & \begin{minipage}[t]{\linewidth}\raggedright
\begin{quote}
Train: 3680 and Test: 921
\end{quote}
\end{minipage} \\
\bottomrule()
\end{longtable}

1

\begin{longtable}[]{@{}
  >{\raggedright\arraybackslash}p{(\columnwidth - 4\tabcolsep) * \real{0.3333}}
  >{\raggedright\arraybackslash}p{(\columnwidth - 4\tabcolsep) * \real{0.3333}}
  >{\raggedright\arraybackslash}p{(\columnwidth - 4\tabcolsep) * \real{0.3333}}@{}}
\toprule()
\multicolumn{2}{@{}>{\raggedright\arraybackslash}p{(\columnwidth - 4\tabcolsep) * \real{0.6667} + 2\tabcolsep}}{%
\begin{minipage}[b]{\linewidth}\raggedright
\begin{quote}
ICS1512
\end{quote}
\end{minipage}} & \begin{minipage}[b]{\linewidth}\raggedright
Machine Learning Algorithms Laboratory
\end{minipage} \\
\midrule()
\endhead
\begin{minipage}[t]{\linewidth}\raggedright
\begin{quote}
\textbf{4}
\end{quote}
\end{minipage} &
\multicolumn{2}{>{\raggedright\arraybackslash}p{(\columnwidth - 4\tabcolsep) * \real{0.6667} + 2\tabcolsep}@{}}{%
\begin{minipage}[t]{\linewidth}\raggedright
\begin{quote}
\textbf{Exploratory Data Analysis}
\end{quote}
\end{minipage}} \\
\bottomrule()
\end{longtable}

\begin{quote}
• Loaded the Spambase dataset and examined its dimensions and data
types.• Generated descriptive statistics using the describe() function.

• Verified that the dataset contains no missing values.

• Visualized the distribution of spam and ham emails.

• Analyzed correlations among numerical features using a heatmap.

• Examined feature distributions using histograms and boxplots to
identify data spread and potential outliers.

\textbf{Class Distribution}
\end{quote}

\includegraphics[width=4.3875in,height=3.13055in]{media/image1.png}

Figure 1: Class Distribution of Emails

\begin{quote}
\textbf{Inference:}\\
The class distribution plot shows that the dataset contains more ham
emails than spam emails. Approximately 60.6\% of the samples belong to
the ham class, while 39.4\% belong to the spam class. Although the
dataset is not perfectly balanced, the class imbalance is moderate and
suitable for training classification models without extensive
resampling.
\end{quote}

2

\begin{longtable}[]{@{}
  >{\raggedright\arraybackslash}p{(\columnwidth - 2\tabcolsep) * \real{0.5000}}
  >{\raggedright\arraybackslash}p{(\columnwidth - 2\tabcolsep) * \real{0.5000}}@{}}
\toprule()
\begin{minipage}[b]{\linewidth}\raggedright
\begin{quote}
ICS1512
\end{quote}
\end{minipage} & \begin{minipage}[b]{\linewidth}\raggedright
Machine Learning Algorithms Laboratory
\end{minipage} \\
\midrule()
\endhead
\bottomrule()
\end{longtable}

\begin{quote}
\textbf{Correlation Matrix}
\end{quote}

\includegraphics[width=4.3875in,height=3.13333in]{media/image2.png}

Figure 2: Correlation Matrix of Features

\begin{quote}
\textbf{Inference:}\\
The correlation matrix illustrates the relationships among the numerical
features in the Spambase dataset. Most features exhibit weak to moderate
correlations, while a few groups of features show stronger positive
correlations. This indicates that the dataset contains diverse
information with limited redundancy, making it suitable for spam
clas-sification.

\textbf{5} \textbf{Data Preprocessing}

• \textbf{Missing Value Handling:}\\
The Spambase dataset contains no missing values; therefore, no
imputation or removal of records was required.

• \textbf{Encoding:}\\
The target variable is already represented as binary labels (Spam/Ham),
so no additional encoding was necessary.

• \textbf{Feature Scaling:}\\
All numerical features were standardized using StandardScaler. Feature
scaling is essential for KNN because it relies on distance calculations.

• \textbf{Train-Test Split:}\\
The dataset was divided into training and testing sets using an 80:20
ratio, resulting in 3680 training samples and 921 testing samples.

• \textbf{Feature Engineering:}\\
No additional feature engineering techniques were applied since the
dataset already contains extracted numerical features.
\end{quote}

3

\begin{longtable}[]{@{}
  >{\raggedright\arraybackslash}p{(\columnwidth - 4\tabcolsep) * \real{0.3333}}
  >{\raggedright\arraybackslash}p{(\columnwidth - 4\tabcolsep) * \real{0.3333}}
  >{\raggedright\arraybackslash}p{(\columnwidth - 4\tabcolsep) * \real{0.3333}}@{}}
\toprule()
\multicolumn{2}{@{}>{\raggedright\arraybackslash}p{(\columnwidth - 4\tabcolsep) * \real{0.6667} + 2\tabcolsep}}{%
\begin{minipage}[b]{\linewidth}\raggedright
\begin{quote}
ICS1512
\end{quote}
\end{minipage}} & \begin{minipage}[b]{\linewidth}\raggedright
Machine Learning Algorithms Laboratory
\end{minipage} \\
\midrule()
\endhead
\begin{minipage}[t]{\linewidth}\raggedright
\begin{quote}
\textbf{6}
\end{quote}
\end{minipage} &
\multicolumn{2}{>{\raggedright\arraybackslash}p{(\columnwidth - 4\tabcolsep) * \real{0.6667} + 2\tabcolsep}@{}}{%
\begin{minipage}[t]{\linewidth}\raggedright
\begin{quote}
\textbf{Machine Learning Algorithm}
\end{quote}
\end{minipage}} \\
\bottomrule()
\end{longtable}

\begin{quote}
• \textbf{Working Principle:}

\textbf{-- Naive Bayes:} Naive Bayes is a probabilistic supervised
learning algorithm based on Bayes' theorem. It assumes that all features
are conditionally inde-pendent given the class label and predicts the
class with the highest posterior probability.

\textbf{-- K-Nearest Neighbors (KNN):} KNN is a non-parametric
supervised learn-ing algorithm that classifies a sample based on the
majority class among its \emph{k} nearest neighbors using a distance
metric such as Euclidean or Manhattan distance.

• \textbf{Mathematical Formulation:}

Bayes' theorem is given by
\end{quote}

\[P(C|X) = \frac{P(X|C) \times P(C)}{P(X)}\]

\begin{quote}
where \emph{P}(\emph{C\textbar X}) is the posterior probability,
\emph{P}(\emph{X\textbar C}) is the likelihood, \emph{P}(\emph{C}) is
the prior probability and \emph{P}(\emph{X}) is the evidence.

The Manhattan distance used in KNN is
\end{quote}

\[d(x, y) = \sum_{i=1}^{n} |x_i - y_i|\]

\begin{quote}
where \emph{x} and \emph{y} are two data points.
\end{quote}

\begin{quote}
• \textbf{Advantages:}

\textbf{--} Naive Bayes is simple, fast and works well for text
classification.

\textbf{--} KNN is easy to implement and can model non-linear decision
boundaries.

\textbf{--} Both algorithms are effective for binary classification
problems such as spam detection.

• \textbf{Limitations:}

\textbf{--} Naive Bayes assumes feature independence, which may not
always hold.

\textbf{--} KNN has high prediction time because it computes distances
to all training samples.

\textbf{--} KNN is sensitive to the choice of \emph{k} and requires
feature scaling.
\end{quote}

\begin{longtable}[]{@{}
  >{\raggedright\arraybackslash}p{(\columnwidth - 4\tabcolsep) * \real{0.3333}}
  >{\raggedright\arraybackslash}p{(\columnwidth - 4\tabcolsep) * \real{0.3333}}
  >{\raggedright\arraybackslash}p{(\columnwidth - 4\tabcolsep) * \real{0.3333}}@{}}
\toprule()
\begin{minipage}[b]{\linewidth}\raggedright
\begin{quote}
\textbf{7}
\end{quote}
\end{minipage} &
\multicolumn{2}{>{\raggedright\arraybackslash}p{(\columnwidth - 4\tabcolsep) * \real{0.6667} + 2\tabcolsep}@{}}{%
\begin{minipage}[b]{\linewidth}\raggedright
\begin{quote}
\textbf{Experimental Setup}
\end{quote}
\end{minipage}} \\
\midrule()
\endhead
\multicolumn{2}{@{}>{\raggedright\arraybackslash}p{(\columnwidth - 4\tabcolsep) * \real{0.6667} + 2\tabcolsep}}{%
\begin{minipage}[t]{\linewidth}\raggedright
\begin{quote}
Python Version\\
Scikit-Learn Version IDE\\
Operating System\\
Hardware
\end{quote}\strut
\end{minipage}} & \begin{minipage}[t]{\linewidth}\raggedright
\begin{quote}
3.13.13\\
1.7.5\\
Jupyter Notebook\\
Fedora 42\\
Intel Core i5 (11th gen), 8GB RAM
\end{quote}\strut
\end{minipage} \\
\bottomrule()
\end{longtable}

4

\begin{longtable}[]{@{}
  >{\raggedright\arraybackslash}p{(\columnwidth - 12\tabcolsep) * \real{0.1429}}
  >{\raggedright\arraybackslash}p{(\columnwidth - 12\tabcolsep) * \real{0.1429}}
  >{\raggedright\arraybackslash}p{(\columnwidth - 12\tabcolsep) * \real{0.1429}}
  >{\raggedright\arraybackslash}p{(\columnwidth - 12\tabcolsep) * \real{0.1429}}
  >{\raggedright\arraybackslash}p{(\columnwidth - 12\tabcolsep) * \real{0.1429}}
  >{\raggedright\arraybackslash}p{(\columnwidth - 12\tabcolsep) * \real{0.1429}}
  >{\raggedright\arraybackslash}p{(\columnwidth - 12\tabcolsep) * \real{0.1429}}@{}}
\toprule()
\multicolumn{2}{@{}>{\raggedright\arraybackslash}p{(\columnwidth - 12\tabcolsep) * \real{0.2857} + 2\tabcolsep}}{%
\begin{minipage}[b]{\linewidth}\raggedright
\begin{quote}
ICS1512
\end{quote}
\end{minipage}} &
\multicolumn{5}{>{\raggedright\arraybackslash}p{(\columnwidth - 12\tabcolsep) * \real{0.7143} + 8\tabcolsep}@{}}{%
\begin{minipage}[b]{\linewidth}\raggedright
Machine Learning Algorithms Laboratory
\end{minipage}} \\
\midrule()
\endhead
\begin{minipage}[t]{\linewidth}\raggedright
\begin{quote}
\textbf{8}
\end{quote}
\end{minipage} &
\multicolumn{6}{>{\raggedright\arraybackslash}p{(\columnwidth - 12\tabcolsep) * \real{0.8571} + 10\tabcolsep}@{}}{%
\begin{minipage}[t]{\linewidth}\raggedright
\begin{quote}
\textbf{Hyperparameter Selection}
\end{quote}
\end{minipage}} \\
\multicolumn{2}{@{}>{\raggedright\arraybackslash}p{(\columnwidth - 12\tabcolsep) * \real{0.2857} + 2\tabcolsep}}{%
\begin{minipage}[t]{\linewidth}\raggedright
\begin{quote}
Parameter
\end{quote}
\end{minipage}} &
\multicolumn{5}{>{\raggedright\arraybackslash}p{(\columnwidth - 12\tabcolsep) * \real{0.7143} + 8\tabcolsep}@{}}{%
\begin{minipage}[t]{\linewidth}\raggedright
\begin{quote}
Value
\end{quote}
\end{minipage}} \\
\multicolumn{2}{@{}>{\raggedright\arraybackslash}p{(\columnwidth - 12\tabcolsep) * \real{0.2857} + 2\tabcolsep}}{%
\begin{minipage}[t]{\linewidth}\raggedright
\begin{quote}
Number of Neighbors (\emph{k}) Distance Metric\\
Search Algorithm\\
Weight Function\\
Cross Validation
\end{quote}\strut
\end{minipage}} &
\multicolumn{5}{>{\raggedright\arraybackslash}p{(\columnwidth - 12\tabcolsep) * \real{0.7143} + 8\tabcolsep}@{}}{%
\begin{minipage}[t]{\linewidth}\raggedright
\begin{quote}
8\\
Manhattan\\
KDTree and BallTree (Compared)\\
Distance\\
5-Fold
\end{quote}\strut
\end{minipage}} \\
\begin{minipage}[t]{\linewidth}\raggedright
\begin{quote}
\textbf{9}
\end{quote}
\end{minipage} &
\multicolumn{6}{>{\raggedright\arraybackslash}p{(\columnwidth - 12\tabcolsep) * \real{0.8571} + 10\tabcolsep}@{}}{%
\begin{minipage}[t]{\linewidth}\raggedright
\begin{quote}
\textbf{Results}
\end{quote}
\end{minipage}} \\
\multicolumn{2}{@{}>{\raggedright\arraybackslash}p{(\columnwidth - 12\tabcolsep) * \real{0.2857} + 2\tabcolsep}}{%
\multirow{2}{*}{Model}} &
\multirow{2}{*}{\begin{minipage}[t]{\linewidth}\raggedright
\begin{quote}
Accuracy
\end{quote}
\end{minipage}} & \multirow{2}{*}{Precision} & &
\multirow{2}{*}{\begin{minipage}[t]{\linewidth}\raggedright
\begin{quote}
F1-Score
\end{quote}
\end{minipage}} &
\multirow{2}{*}{\begin{minipage}[t]{\linewidth}\raggedright
\begin{quote}
ROC-AUC
\end{quote}
\end{minipage}} \\
& & & & \begin{minipage}[t]{\linewidth}\raggedright
\begin{quote}
Recall
\end{quote}
\end{minipage} \\
\multicolumn{2}{@{}>{\raggedright\arraybackslash}p{(\columnwidth - 12\tabcolsep) * \real{0.2857} + 2\tabcolsep}}{%
\begin{minipage}[t]{\linewidth}\raggedright
\begin{quote}
Gaussian NB\\
Multinomial NB\\
Bernoulli NB\\
Optimized KNN
\end{quote}\strut
\end{minipage}} & \begin{minipage}[t]{\linewidth}\raggedright
\begin{quote}
0.8328\\ 0.8936\\ 0.8795\\ 0.9142
\end{quote}
\end{minipage} & \begin{minipage}[t]{\linewidth}\raggedright
\begin{quote}
0.7146\\ 0.9431\\ 0.8684\\ 0.9356
\end{quote}
\end{minipage} & \begin{minipage}[t]{\linewidth}\raggedright
\begin{quote}
0.9587\\ 0.7769\\ 0.8182\\ 0.8402
\end{quote}
\end{minipage} & \begin{minipage}[t]{\linewidth}\raggedright
\begin{quote}
0.8188\\ 0.8520\\ 0.8426\\ 0.8853
\end{quote}
\end{minipage} & \begin{minipage}[t]{\linewidth}\raggedright
\begin{quote}
0.9229\\
0.9625\\
0.9457\\
0.9563
\end{quote}\strut
\end{minipage} \\
\bottomrule()
\end{longtable}

Table 2: Performance Comparison of Classification Models

\begin{quote}
\textbf{10 Result Visualization}

\textbf{Confusion matrix optimized KNN}
\end{quote}

\includegraphics[width=4.3875in,height=4.00278in]{media/image3.png}

Figure 3: Confusion Matrix of Optimized KNN

5

\begin{longtable}[]{@{}
  >{\raggedright\arraybackslash}p{(\columnwidth - 2\tabcolsep) * \real{0.5000}}
  >{\raggedright\arraybackslash}p{(\columnwidth - 2\tabcolsep) * \real{0.5000}}@{}}
\toprule()
\begin{minipage}[b]{\linewidth}\raggedright
\begin{quote}
ICS1512
\end{quote}
\end{minipage} & \begin{minipage}[b]{\linewidth}\raggedright
Machine Learning Algorithms Laboratory
\end{minipage} \\
\midrule()
\endhead
\bottomrule()
\end{longtable}

\begin{quote}
\textbf{Inference:} The confusion matrix indicates that the optimized
KNN classifier correctly classified most spam and ham emails with only a
small number of misclassifications. The high number of true positives
and true negatives demonstrates the effectiveness of the classifier. The
low false positive and false negative counts indicate reliable
prediction performance.

\textbf{ROC Curve}
\end{quote}

\includegraphics[width=4.3875in,height=3.47222in]{media/image4.png}

Figure 4: ROC Curve of the Optimized KNN Classifier

\begin{quote}
\textbf{Inference:}\\
The ROC curve illustrates the trade-off between the True Positive Rate
(TPR) and False Positive Rate (FPR) at different classification
thresholds. The curve lies close to the upper-left corner, indicating
that the classifier effectively distinguishes spam emails from ham
emails. A high Area Under the Curve (AUC) value demonstrates excellent
classification performance and confirms the robustness of the optimized
KNN model.

\textbf{11 Discussion}

The experiment compared three Naive Bayes variants with the KNN
algorithm for email spam classification. Among the Naive Bayes
classifiers, Multinomial Naive Bayes achieved the highest accuracy.
Hyperparameter tuning using GridSearchCV further improved the
performance of KNN, resulting in the highest overall accuracy. The
comparison between KDTree and BallTree showed similar classification
accuracy with differences only in execution time. K-Fold Cross
Validation confirmed that the models produced consistent performance
across different data splits.
\end{quote}

6

\begin{longtable}[]{@{}
  >{\raggedright\arraybackslash}p{(\columnwidth - 4\tabcolsep) * \real{0.3333}}
  >{\raggedright\arraybackslash}p{(\columnwidth - 4\tabcolsep) * \real{0.3333}}
  >{\raggedright\arraybackslash}p{(\columnwidth - 4\tabcolsep) * \real{0.3333}}@{}}
\toprule()
\multicolumn{2}{@{}>{\raggedright\arraybackslash}p{(\columnwidth - 4\tabcolsep) * \real{0.6667} + 2\tabcolsep}}{%
\begin{minipage}[b]{\linewidth}\raggedright
\begin{quote}
ICS1512
\end{quote}
\end{minipage}} & \begin{minipage}[b]{\linewidth}\raggedright
\begin{quote}
Machine Learning Algorithms Laboratory
\end{quote}
\end{minipage} \\
\midrule()
\endhead
\begin{minipage}[t]{\linewidth}\raggedright
\begin{quote}
\textbf{12}
\end{quote}
\end{minipage} & \begin{minipage}[t]{\linewidth}\raggedright
\begin{quote}
\textbf{Conclusion}
\end{quote}
\end{minipage} & \\
\bottomrule()
\end{longtable}

\begin{quote}
The experiment successfully implemented Email Spam/Ham classification
using Naive Bayes and KNN algorithms. Among the Naive Bayes variants,
Multinomial Naive Bayes achieved the best performance. Hyperparameter
optimization further improved KNN, making it the best-performing model
with an accuracy of approximately 91.42\%. The study demonstrated the
effectiveness of hyperparameter tuning and cross-validation in improving
model performance.

\textbf{13 References}

1. A. Géron, \emph{Hands-On Machine Learning with Scikit-Learn, Keras \&
TensorFlow}, 3rd ed., O'Reilly Media, 2022.

2. T. Hastie, R. Tibshirani and J. Friedman, \emph{The Elements of
Statistical Learning}. Springer, 2009.

3. Scikit-Learn Documentation. Available: \url{https://scikit-learn.org}

4. Kaggle Spambase Dataset. Available:
\url{https://www.kaggle.com/datasets/somesh24/spambase}
\end{quote}

7

\end{document}
